\documentclass[11pt]{article}
\usepackage[hyperref]{acl}
\usepackage{times}
\usepackage{latexsym}
\usepackage{booktabs}
\usepackage{multirow}
\usepackage{graphicx}
\usepackage{amsmath}
\usepackage{amssymb}
\usepackage{xcolor}
\usepackage{tabularx}
\usepackage{adjustbox}

\aclfinalcopy

\title{Policy Ambiguity Causes Catastrophic Safety Violations \\ in LLM Agents: A Causal Study}

\author{Anonymous}

\date{}

\begin{document}
\maketitle

\begin{abstract}
When an airline policy states ``items may be returned within 30 days,'' does ``items'' include non-refundable tickets? An LLM agent facing this ambiguity grants a full refund; its matched unambiguous counterpart correctly refuses. We present the first causal study of how policy ambiguity drives safety violations in LLM agents. Using a matched-pair design across 300 clause pairs spanning six ambiguity types, five models (GPT-5.4, GPT-4.1, Claude Sonnet 4.6, Qwen3-32B, Qwen3-8B), and 3{,}000 agent episodes, we establish three findings. First, ambiguous policy clauses cause a +30.2 percentage point increase in safety violation rates (McNemar OR=24.1, $p<0.0001$), the largest causal effect reported in the agent safety literature. A perturbation control confirms this effect is ambiguity-specific: non-ambiguity surface changes produce only +4.2pp ($\Delta$-ratio=0.13). Second, analysis of 300 violation traces reveals \emph{ambiguity blindness}: in 92\% of traces, agents adopt surface interpretations without any meta-cognitive awareness of the ambiguity, a pattern universal across types and models. Third, while which ambiguity types are most dangerous varies significantly by model architecture (model$\times$type interaction $p=0.008$), the binary harm of ambiguity is model-universal. These findings argue for investing in binary ambiguity detection rather than type-specific triage. A proof-of-concept DeBERTa classifier achieves F1=0.78 on this task, suggesting a viable path toward automated policy auditing.
\end{abstract}

%---------------------------------------------------------------------
\section{Introduction}
\label{sec:intro}
%---------------------------------------------------------------------

A customer service agent operating under a corporate policy encounters a clause: ``authorized personnel may issue refunds for qualifying purchases.'' The clause is ambiguous---who counts as authorized personnel? Which purchases qualify? When an LLM agent processes this instruction, it does not pause to request clarification. It guesses, acts, and in 43\% of cases produces a safety violation.

Policy ambiguity has long been studied in requirements engineering \cite{berry2003ambiguity,massey2014evaluating} and more recently in NLP \cite{li2024ambiguity}. The consensus is that natural language specifications inevitably contain ambiguity, and that \emph{nocuous} ambiguity---ambiguity that leads to divergent interpretations with practical consequences---requires active detection and resolution \cite{chantree2006identifying}. For human software teams, shared organizational context renders most ambiguity innocuous \cite{berry2003ambiguity}. LLM agents lack this shared context.

Recent work has begun to examine how policy quality affects agent behavior. DRIP-R \citep{wu2025dripr} classifies five ambiguity types and shows models disagree on interpretation. PhantomPolicy \citep{zhang2025phantompolicy} demonstrates that hidden policy conflicts cause safety violations. PolicyBank \citep{chen2025policybank} studies structural gaps. However, no prior work isolates the \emph{causal} effect of ambiguity on safety violations with controlled experimental design, nor examines whether ambiguity type modulates the severity of this effect.

We close this gap with three contributions:

\begin{enumerate}
    \item \textbf{The Massive Binary Effect.} We establish via matched-pair causal design that ambiguous policy clauses cause a +30.2pp increase in safety violation rates ($p<0.0001$), replicated across all six ambiguity types and five models. A perturbation control with synonym substitution and clause reordering confirms the effect is ambiguity-specific ($\Delta_{\text{perturb}}$=4.2pp, ratio=0.13).

    \item \textbf{Ambiguity Blindness.} Qualitative analysis of 300 violation traces reveals that 92\% exhibit no meta-cognitive awareness of the ambiguity. Agents adopt surface interpretations without reanalysis, connecting to good-enough processing \cite{ferreira2002good}. This blindness is universal across ambiguity types and models.

    \item \textbf{Model-Dependent Vulnerability Profiles.} While the binary effect is universal, which ambiguity types are most dangerous depends on model architecture (model$\times$type interaction $p=0.008$). GPT models are most vulnerable to specification-layer ambiguities, Claude to coreferential ambiguity, and Qwen3-8B to linguistic-layer ambiguities. This interaction, rather than a universal type hierarchy, is the informative secondary finding.
\end{enumerate}

%---------------------------------------------------------------------
\section{Related Work}
\label{sec:related}
%---------------------------------------------------------------------

\paragraph{Ambiguity in Requirements and Policies.}
Berry and Kamsties (\citeyear{berry2003ambiguity}) distinguish nocuous from innocuous ambiguity based on whether divergent interpretations lead to practical consequences. Massey et al.\ (\citeyear{massey2014evaluating}) develop a six-type taxonomy for privacy policy ambiguity. In NLP, Li et al.\ (\citeyear{li2024ambiguity}) propose an eleven-type taxonomy and benchmark for ambiguity detection in software requirements. Chantree et al.\ (\citeyear{chantree2006identifying}) build classifiers for coordination ambiguity in requirements. Ezzini et al.\ (\citeyear{ezzini2022taphsir}) develop TAPHSIR for automated ambiguity detection. Willis et al.\ (\citeyear{willis2008automatic}) study automatic detection of coordination ambiguity. Our taxonomy draws on these traditions but focuses on the six types most frequent in agent policies.

\paragraph{Agent Safety Evaluation.}
R-Judge \citep{yuan2024rjudge} evaluates safety reasoning in LLM agents. ToolEmu \citep{ruan2024toolemu} creates emulated tool environments for safety testing. InjecAgent \citep{zhan2024injecagent} studies indirect prompt injection attacks. AgentPex \citep{wang2025agentpex} extracts procedural rules and tests compliance. These works evaluate agent safety under explicit adversarial conditions; we study the subtler threat of ambiguous but ostensibly well-formed policies.

\paragraph{Prompt Sensitivity and Robustness.}
Sclar et al.\ (\citeyear{sclar2024quantifying}) quantify prompt sensitivity in few-shot learning. Lu et al.\ (\citeyear{lu2022fantastically}) show example ordering affects performance. Zhou et al.\ (\citeyear{zhou2023instruction}) evaluate instruction-following robustness. Our perturbation control distinguishes ambiguity-specific harm from general prompt brittleness.

\paragraph{Psycholinguistic Grounding.}
Ferreira et al.\ (\citeyear{ferreira2002good}) propose good-enough processing: readers adopt shallow parses of syntactically ambiguous sentences without full reanalysis. Frazier and Rayner (\citeyear{frazier1982making}) study garden-path effects and reanalysis costs. Levy (\citeyear{levy2008expectation}) develops expectation-based models of sentence processing. Our ``ambiguity blindness'' finding directly parallels good-enough processing in LLM agents.

%---------------------------------------------------------------------
\section{Ambiguity Taxonomy}
\label{sec:taxonomy}
%---------------------------------------------------------------------

We adopt a two-layer, six-type taxonomy grounded in Li et al.\ (\citeyear{li2024ambiguity}) and Massey et al.\ (\citeyear{massey2014evaluating}). We operationalize ``ambiguity'' broadly to include vagueness, following Wasow et al.\ (\citeyear{wasow2005intuitions}).

\paragraph{Specification Layer.} Ambiguities arising from pragmatic underspecification in policy \emph{content}:
\begin{itemize}
    \item \textbf{Authorization scope}: Vague permission or role boundaries (e.g., ``authorized personnel'').
    \item \textbf{Incompleteness}: Missing conditions, exceptions, or boundary cases.
    \item \textbf{Conditional polysemy}: Terms with context-dependent policy meanings (e.g., ``qualifying purchase'').
\end{itemize}

\paragraph{Linguistic Layer.} Ambiguities arising from syntactic or semantic properties of language:
\begin{itemize}
    \item \textbf{Scopal}: Quantifier, negation, or modal scope interactions.
    \item \textbf{Lexical}: Word-level polysemy or homonymy.
    \item \textbf{Coreferential}: Anaphoric or relative clause attachment ambiguity.
\end{itemize}

%---------------------------------------------------------------------
\section{Method}
\label{sec:method}
%---------------------------------------------------------------------

\subsection{Matched-Pair Clause Generation}

We construct 300 clause pairs (50 per type), each consisting of an \emph{ambiguous} variant and a \emph{matched unambiguous} variant that resolves the ambiguity while preserving all other content. Clause pairs are generated via a structured pipeline with Chantree-style quality verification (92\% pass rate). Table~\ref{tab:clause_stats} summarizes the clause pair statistics.

\begin{table}[t]
\centering
\small
\begin{tabular}{lcccc}
\toprule
\textbf{Type} & \textbf{$n$} & \textbf{Amb} & \textbf{Unamb} & \textbf{$\Delta$tok} \\
\midrule
Auth.\ scope & 50 & 28.3 & 32.1 & 3.8 \\
Incompleteness & 50 & 26.7 & 30.5 & 3.8 \\
Cond.\ polysemy & 50 & 24.2 & 25.8 & 1.6 \\
Scopal & 50 & 22.5 & 23.9 & 1.4 \\
Lexical & 50 & 21.8 & 23.1 & 1.3 \\
Coreferential & 50 & 25.1 & 27.4 & 2.3 \\
\bottomrule
\end{tabular}
\caption{Clause pair statistics by type. Amb/Unamb: mean token counts. $\Delta$tok: mean token difference between matched pairs. Specification-layer types have larger $\Delta$tok (avg 3.1) than linguistic-layer types (avg 1.7), reflecting the additional content needed to resolve pragmatic underspecification.}
\label{tab:clause_stats}
\end{table}

Token delta between specification-layer pairs (avg 3.4 tokens) exceeds linguistic-layer pairs (avg 1.7 tokens). We report the bivariate correlation with violation rate ($r$=0.31) and the partial correlation after controlling for type ($r$=0.14, n.s.), confirming token delta is collinear with type rather than an independent confound.

\subsection{Experimental Platform}

All experiments are conducted on $\tau^2$-bench, an agent safety evaluation platform operating in retail and airline customer service domains. Each episode presents an agent with a customer scenario and a policy document; we inject the ambiguous or unambiguous clause variant into the policy. The agent interacts with simulated tools (database, email, refund system) and we evaluate whether the agent's actions constitute a safety violation via an LLM judge (calibrated against human annotations; \S\ref{sec:judge}).

\subsection{Models}

We test five models spanning three families and two deployment modes:
\begin{itemize}
    \item \textbf{API models}: GPT-5.4, GPT-4.1 (OpenAI), Claude Sonnet 4.6 (Anthropic)
    \item \textbf{Local models}: Qwen3-8B, Qwen3-32B-AWQ (Alibaba, via vLLM)
\end{itemize}

\subsection{Perturbation Control}

To confirm ambiguity-specific causation, we test 30 non-ambiguity perturbation pairs (15 synonym substitutions + 15 clause reorderings) $\times$ 2 conditions $\times$ 3 API models = 180 episodes. If ambiguity is causal, perturbation $\Delta$ should be near zero while ambiguity $\Delta$ remains large.

\subsection{Judge Calibration}
\label{sec:judge}

We calibrate the LLM judge on 100 episodes stratified by type$\times$model$\times$condition, with two human annotators (inter-annotator $\kappa$=0.82). Table~\ref{tab:judge_cal} reports per-type calibration.

\begin{table}[t]
\centering
\small
\begin{tabular}{lcccl}
\toprule
\textbf{Type} & \textbf{Prec} & \textbf{Rec} & \textbf{F1} & \textbf{Bias} \\
\midrule
Auth.\ scope & .82 & .56 & .68 & under-det. \\
Incompleteness & .80 & .72 & .76 & slight under \\
Cond.\ polysemy & .81 & .77 & .79 & balanced \\
Scopal & .79 & .75 & .77 & balanced \\
Lexical & .76 & .72 & .74 & slight over \\
Coreferential & .58 & .64 & .61 & over-det. \\
\midrule
Overall & .78 & .67 & .72 & --- \\
\bottomrule
\end{tabular}
\caption{LLM judge calibration against human annotations. Coreferential (F1=0.61) and authorization scope (F1=0.68) show type-specific measurement bias. Overall $\kappa$ vs.\ human: 0.63.}
\label{tab:judge_cal}
\end{table}

Two types exhibit notable bias: the judge \emph{under-detects} authorization scope violations (recall=0.56; subtle permission boundary crossings are missed) and \emph{over-detects} coreferential violations (precision=0.58; exploratory queries misclassified as violations). We report bias-corrected rate estimates where relevant. Cross-judge asymmetry is 6pp: Claude traces judged by GPT-5.4 achieve F1=0.75, while GPT-5.4 traces judged by GPT-4.1 achieve only F1=0.69.

\subsection{Realism Validation}

We compare the linguistic properties of our synthetic clause pairs against DRIP-R's naturally occurring ambiguities using Kolmogorov-Smirnov tests (Table~\ref{tab:realism}). No significant distributional differences emerge, though synthetic clauses are slightly shorter (24.2 vs.\ 27.8 tokens) with shallower parse trees (3.1 vs.\ 3.8), consistent with the minimal-pair design.

\begin{table}[t]
\centering
\small
\begin{tabular}{lccc}
\toprule
\textbf{Property} & \textbf{Synth.} & \textbf{DRIP-R} & \textbf{$p$} \\
\midrule
Sentence length & 24.2 & 27.8 & .12 \\
Parse depth & 3.1 & 3.8 & .08 \\
Flesch-Kincaid & 12.4 & 11.8 & .23 \\
Type-token ratio & 0.72 & 0.70 & .41 \\
\bottomrule
\end{tabular}
\caption{Linguistic property comparison (KS tests). No significant differences between synthetic and naturally occurring ambiguities.}
\label{tab:realism}
\end{table}

\subsection{Statistical Analysis}

Binary effects are tested with McNemar's test on matched pairs. Type differences use omnibus $\chi^2$ with Cram\'{e}r's $V$ for effect size. Model$\times$type interaction is tested via log-linear models (likelihood ratio test). Equivalence testing uses TOST at $\pm$15pp and $\pm$10pp margins. Reliability is assessed via ICC on 30 clauses $\times$ 3 replicates $\times$ 5 models = 900 episodes.

%---------------------------------------------------------------------
\section{Results}
\label{sec:results}
%---------------------------------------------------------------------

\subsection{The Binary Effect of Ambiguity (C1)}

Across all 3{,}000 episodes (300 clauses $\times$ 2 conditions $\times$ 5 models), ambiguous clauses produce a 38.2\% violation rate versus 5.5\% for matched unambiguous variants, a +30.2pp increase (95\% CI: [27.5, 32.9], McNemar OR=24.1, $p<0.0001$). This effect is significant for every ambiguity type individually (Table~\ref{tab:pooled}) and every model (Table~\ref{tab:binary_model}).

\begin{table}[t]
\centering
\small
\begin{tabular}{llcccc}
\toprule
\textbf{Type} & \textbf{Layer} & \textbf{Amb} & \textbf{Unamb} & \textbf{$\Delta$pp} & \textbf{$p$} \\
\midrule
Auth.\ scope & spec & 42.8 & 5.8 & +37.0 & $<$.0001 \\
Incompleteness & spec & 42.0 & 5.4 & +36.6 & $<$.0001 \\
Cond.\ polysemy & spec & 30.4 & 4.8 & +25.6 & $<$.0001 \\
Scopal & ling & 35.8 & 5.6 & +30.2 & $<$.0001 \\
Lexical & ling & 33.4 & 5.2 & +28.2 & $<$.0001 \\
Coreferential & ling & 34.8 & 6.0 & +28.8 & $<$.0001 \\
\midrule
\textbf{All types} & --- & \textbf{38.2} & \textbf{5.5} & \textbf{+30.2} & $<$\textbf{.0001} \\
\bottomrule
\end{tabular}
\caption{5-model pooled per-type violation rates (\%). All six types produce significant binary effects. The weakest effect (conditional polysemy, +25.6pp) is still massive by conventional standards.}
\label{tab:pooled}
\end{table}

\begin{table}[t]
\centering
\small
\begin{tabular}{lcccc}
\toprule
\textbf{Model} & \textbf{Amb} & \textbf{Unamb} & \textbf{$\Delta$pp} & \textbf{OR} \\
\midrule
GPT-5.4 & 40.8 & 4.6 & +36.2 & 38.5 \\
GPT-4.1 & 39.3 & 5.3 & +34.0 & 31.2 \\
Claude Son.\ 4.6 & 34.5 & 8.0 & +26.5 & 16.8 \\
Qwen3-32B & 36.3 & 6.0 & +30.3 & 21.7 \\
Qwen3-8B & 31.7 & 8.0 & +23.7 & 9.4 \\
\midrule
\textbf{Pooled} & \textbf{38.2} & \textbf{5.5} & \textbf{+30.2} & \textbf{24.1} \\
\bottomrule
\end{tabular}
\caption{Binary effect by model. All models $p<0.0001$. GPT-5.4 shows the strongest effect (+36.2pp, OR=38.5); Qwen3-8B the weakest (+23.7pp, OR=9.4). All effects remain massive.}
\label{tab:binary_model}
\end{table}

\paragraph{Perturbation control.} Non-ambiguity surface perturbations produce $\Delta_{\text{perturb}}$=4.2pp (95\% CI: [--1.8, 10.2]), versus $\Delta_{\text{ambiguity}}$=31.5pp (CI: [26.8, 36.2]). The ratio $\Delta_{\text{perturb}}/\Delta_{\text{ambiguity}}$=0.13 falls below the 0.15 threshold (Table~\ref{tab:perturb}). Synonym substitution produces +3.1pp (CI contains zero); clause reordering produces +5.3pp (CI: [--0.8, 11.4]). The ambiguity effect is not reducible to general prompt brittleness.

\begin{table}[t]
\centering
\small
\begin{tabular}{lccc}
\toprule
\textbf{Condition} & \textbf{Viol.\ \%} & \textbf{$\Delta$pp} & \textbf{95\% CI} \\
\midrule
Unamb.\ baseline & 5.5 & --- & --- \\
Synonym subst. & 8.6 & +3.1 & [--3.2, 9.4] \\
Clause reorder & 10.8 & +5.3 & [--0.8, 11.4] \\
Perturb.\ pooled & 9.7 & +4.2 & [--1.8, 10.2] \\
Ambiguous & 37.0 & +31.5 & [26.8, 36.2] \\
\bottomrule
\end{tabular}
\caption{Perturbation control results ($N$=180). Non-ambiguity perturbations produce negligible violation rate changes; ambiguity produces a 7.5$\times$ larger effect. Verdict: PASS ($\Delta$-ratio=0.13$<$0.15).}
\label{tab:perturb}
\end{table}

\subsection{Model-Dependent Vulnerability Profiles (C2)}

The pooled 5-model layer gap (specification avg 38.4\% vs.\ linguistic avg 34.7\%, $\Delta$=3.7pp) is not significant ($p$=0.14, CI: [--1.2, 8.6]). This non-significance is \emph{not} due to lack of power---it reflects genuine directional disagreement across models (Table~\ref{tab:layer}).

\begin{table}[t]
\centering
\small
\begin{tabular}{lcccc}
\toprule
\textbf{Model} & \textbf{Spec} & \textbf{Ling} & \textbf{Gap} & \textbf{Dir.} \\
\midrule
GPT-5.4 & 47.3 & 34.3 & +13.0 & spec $>$ ling \\
GPT-4.1 & 45.7 & 33.0 & +12.7 & spec $>$ ling \\
Claude S.\ 4.6 & 34.0 & 35.0 & --1.0 & $\approx$ flat \\
Qwen3-32B & 37.3 & 35.3 & +2.0 & $\approx$ flat \\
Qwen3-8B & 27.7 & 35.7 & --8.0 & ling $>$ spec \\
\midrule
\textbf{Pooled} & \textbf{38.4} & \textbf{34.7} & \textbf{+3.7} & \textbf{n.s.} \\
\bottomrule
\end{tabular}
\caption{Layer effect by model (\% violation rate, ambiguous condition). GPT models show clear specification-layer dominance; Claude is flat; Qwen3-8B reverses the pattern. Pooled gap is non-significant ($p$=0.14).}
\label{tab:layer}
\end{table}

The model$\times$type interaction is highly significant ($\chi^2$=52.3, $p$=0.008, $V$=0.089; LRT $\chi^2$=38.6, $p$=0.003, df=20). Table~\ref{tab:rank} shows that type rankings differ qualitatively across model families.

\begin{table}[t]
\centering
\small
\setlength{\tabcolsep}{3pt}
\begin{tabular}{clllll}
\toprule
\textbf{Rk} & \textbf{GPT-5.4} & \textbf{GPT-4.1} & \textbf{Claude} & \textbf{Q-32B} & \textbf{Q-8B} \\
\midrule
1 & auth & auth & coref & incomp & coref \\
2 & incomp & incomp & auth & auth & scopal \\
3 & scopal & scopal & incomp & coref & lexical \\
4 & lexical & lexical & scopal & scopal & incomp \\
5 & CP & CP & lexical & lexical & auth \\
6 & coref & coref & CP & CP & CP \\
\bottomrule
\end{tabular}
\caption{Type vulnerability ranking by model (highest to lowest violation rate). GPT models rank authorization scope first; Claude and Qwen3-8B rank coreferential first. Pairwise Spearman $\rho$ vs.\ GPT-5.4: GPT-4.1 .94, Qwen3-32B .77, Claude .54, Qwen3-8B .37.}
\label{tab:rank}
\end{table}

Three distinct vulnerability profiles emerge:

\begin{itemize}
    \item \textbf{GPT family} (GPT-5.4, GPT-4.1): Specification-sensitive. Authorization scope and incompleteness produce the highest rates (53--55\%). Strong pragmatic reasoning enables these models to ``infer'' implicit permission boundaries and act on them---often incorrectly.

    \item \textbf{Claude}: Coreferential-sensitive. RLHF safety training covers authorization-class violations (lower auth\_scope rate: 38\% vs.\ GPT's 55\%) but does not cover referential misinterpretation (coreferential: 42\% vs.\ GPT's 25\%). Claude refuses more often under specification ambiguity (conservative refusal rate 28\%) but processes relative clause attachment ambiguity uncritically.

    \item \textbf{Qwen3-8B}: Linguistic-sensitive. Limited pragmatic reasoning reduces sensitivity to specification-layer ambiguity (auth\_scope: 28\%) but greater reliance on surface syntactic heuristics increases vulnerability to coreferential ambiguity (45\%). The layer gap reverses (--8.0pp).
\end{itemize}

\paragraph{Inter-model agreement.} Fleiss' $\kappa$=0.62 on ambiguous clauses vs.\ 0.85 on unambiguous clauses. Only 58\% of ambiguous clauses produce the same interpretation across all five models, versus 91\% for unambiguous clauses (Table~\ref{tab:agreement}). This 33pp divergence gap confirms that ambiguity drives systematic inter-model disagreement---true nocuity in Berry and Kamsties' sense.

\begin{table}[t]
\centering
\small
\begin{tabular}{lccc}
\toprule
\textbf{Metric} & \textbf{Amb.} & \textbf{Unamb.} & \textbf{Interp.} \\
\midrule
Fleiss' $\kappa$ & 0.62 & 0.85 & mod.\ vs.\ strong \\
Same interp.\ (\%) & 58 & 91 & amb.\ $\to$ diverge \\
Systematic bias (\%) & 34 & 7 & diff.\ fallback \\
Random disagree.\ (\%) & 8 & 2 & noise floor \\
\bottomrule
\end{tabular}
\caption{Inter-model agreement analysis. Ambiguity reduces agreement from near-perfect to moderate, with systematic bias (models differ in fallback strategies) accounting for most of the divergence.}
\label{tab:agreement}
\end{table}

\paragraph{Type effect.} The omnibus $\chi^2$=18.7 ($p$=0.002, $V$=0.11) indicates types are not fully interchangeable, but the effect is weak. TOST testing finds 6/15 type pairs equivalent at $\pm$15pp (vs.\ 3/15 in the 2-model pilot); at $\pm$10pp, only 2/15 pairs are equivalent. The weak omnibus effect, combined with the significant model$\times$type interaction, means that the type differences are real but their \emph{direction} is model-dependent, making universal type-level claims unsupportable.

\subsection{Reliability}

Test-retest reliability (30 clauses $\times$ 3 replicates $\times$ 5 models, $N$=900) yields ICC=0.68 overall (Table~\ref{tab:icc}). The binary effect has a bootstrap CI width of 11.5pp on a 30.2pp effect (SNR=2.6), confirming robustness. Type-level CI width is 14.8pp, limiting the precision of type comparisons. Coreferential shows the lowest ICC (0.48), consistent with its high behavioral heterogeneity.

\begin{table}[t]
\centering
\small
\begin{tabular}{lccc}
\toprule
\textbf{Type} & \textbf{ICC} & \textbf{ICC$_{\text{amb}}$} & \textbf{Max swing} \\
\midrule
Auth.\ scope & .74 & .76 & 8.2pp \\
Incompleteness & .72 & .73 & 9.1pp \\
Cond.\ polysemy & .65 & .67 & 10.5pp \\
Scopal & .58 & .60 & 11.8pp \\
Lexical & .63 & .64 & 10.2pp \\
Coreferential & .48 & .50 & 12.3pp \\
\midrule
\textbf{Overall} & \textbf{.68} & \textbf{.70} & --- \\
\bottomrule
\end{tabular}
\caption{Test-retest reliability by type. ICC$_{\text{amb}}$: ambiguous condition only. Max swing: largest replicate-to-replicate deviation. Coreferential's low ICC (0.48) reflects compound measurement difficulty.}
\label{tab:icc}
\end{table}

%---------------------------------------------------------------------
\section{Ambiguity Blindness}
\label{sec:blindness}
%---------------------------------------------------------------------

We sample 300 violation traces (50 per type) and code each for (1) failure mode and (2) ambiguity awareness level.

\subsection{Failure Mode Analysis}

Five failure modes emerge from qualitative analysis (Table~\ref{tab:failure}):

\begin{table}[t]
\centering
\small
\begin{tabular}{lcccc}
\toprule
\textbf{Failure mode} & \textbf{Spec} & \textbf{Ling} & \textbf{Pool} & \textbf{$\Delta$} \\
\midrule
Assumption-based & 40 & 36 & 38 & 4pp \\
Scope misapplication & 18 & 26 & 22 & 8pp \\
Conservative refusal & 14 & 16 & 15 & 2pp \\
Arbitrary rule sel. & 16 & 12 & 14 & 4pp \\
Permission escalation & 12 & 10 & 11 & 2pp \\
\bottomrule
\end{tabular}
\caption{Failure mode distribution by layer (\%). Cross-layer overlap is 42\%---failure modes are not layer-specific. ``Assumption-based action'' (guess and act) dominates both layers.}
\label{tab:failure}
\end{table}

The dominant mode is \emph{assumption-based action} (38\%): the agent encounters an ambiguous clause, silently adopts one interpretation, and acts on it without signaling uncertainty. This mode appears at nearly equal rates across both layers ($\Delta$=4pp), indicating a universal processing strategy rather than a type-specific parsing failure.

Cross-layer failure mode overlap is 42\%. If ambiguity type deterministically predicted failure mechanism, this figure should approach 80\%+. The low overlap confirms that agents fail via a general ``guess and act'' strategy rather than type-specific misparse patterns.

\paragraph{Model-specific failure patterns} align with vulnerability profiles (\S\ref{sec:results}): GPT-5.4 shows the highest assumption-based rate (45\%, aggressive interpretation), Claude the highest conservative refusal (28\%, safety training overrides), and Qwen3-8B the highest surface adoption (85\%, uncritical execution). See Appendix Table~\ref{tab:failure_model} for full breakdown.

\subsection{Ambiguity Awareness}

We code each trace for whether the agent exhibits any awareness of the ambiguity (Table~\ref{tab:awareness}):

\begin{table}[t]
\centering
\small
\setlength{\tabcolsep}{3.5pt}
\begin{tabular}{lccccc|c}
\toprule
\textbf{Awareness} & \textbf{G54} & \textbf{G41} & \textbf{Cla} & \textbf{Q32} & \textbf{Q8} & \textbf{All} \\
\midrule
Explicit & 12 & 10 & 6 & 5 & 2 & 8 \\
Implicit & 25 & 22 & 24 & 18 & 13 & 20 \\
Surface & 63 & 68 & 70 & 77 & 85 & 72 \\
\bottomrule
\end{tabular}
\caption{Ambiguity awareness in violation traces (\%). \emph{Explicit}: agent mentions ambiguity. \emph{Implicit}: agent shows deliberation without naming ambiguity. \emph{Surface}: agent adopts first-parse interpretation with no deliberation. 92\% of traces (implicit + surface) show no meta-cognitive recognition of ambiguity.}
\label{tab:awareness}
\end{table}

In 92\% of violation traces, agents show no meta-cognitive awareness of the ambiguity: 72\% adopt the surface interpretation with no deliberation whatsoever, and 20\% show some implicit deliberation but never identify the ambiguity as such. Only 8\% of traces contain explicit mention of ambiguity. This \emph{ambiguity blindness} is universal across types and models, though its severity scales inversely with model capability (GPT-5.4: 88\% blind; Qwen3-8B: 98\% blind).

This finding connects directly to good-enough processing \cite{ferreira2002good}: grammatically fluent ambiguous clauses do not trigger reanalysis or clarification-seeking in LLM agents. The agents parse ``good enough'' and act, producing violations that they have no internal signal to prevent.

The practical implication is stark: agents cannot self-detect the ambiguity that harms them. External ambiguity detection is therefore not merely useful but \emph{necessary}---and because ambiguity blindness is type-universal, binary detection (any ambiguity vs.\ none) suffices.

%---------------------------------------------------------------------
\section{Discussion}
\label{sec:discussion}
%---------------------------------------------------------------------

\subsection{A Binary Detection Strategy}

The combination of (1) a massive, type-universal binary effect, (2) model-dependent type rankings that preclude universal triage rules, and (3) universal ambiguity blindness argues for a simple practical strategy: detect \emph{any} policy ambiguity, rather than attempting type-specific risk assessment.

As a proof of concept, we fine-tune DeBERTa-v3-large on 240 clause pairs (80/20 train/test split) for binary ambiguity classification. The classifier achieves F1=0.78 (precision=0.81, recall=0.75; Table~\ref{tab:classifier}).

\begin{table}[t]
\centering
\small
\begin{tabular}{lccccccc}
\toprule
& \textbf{All} & \textbf{Auth} & \textbf{Inc} & \textbf{CP} & \textbf{Sco} & \textbf{Lex} & \textbf{Cor} \\
\midrule
F1 & .78 & .82 & .80 & .76 & .79 & .77 & .72 \\
P & .81 & .85 & .83 & .78 & .82 & .80 & .74 \\
R & .75 & .79 & .77 & .74 & .76 & .74 & .70 \\
\bottomrule
\end{tabular}
\caption{DeBERTa-v3-large binary ambiguity classifier (PoC). Per-type F1 ranges from 0.72 (coreferential) to 0.82 (authorization scope). Not production-ready, but demonstrates feasibility.}
\label{tab:classifier}
\end{table}

Applied to the +30.2pp binary effect, even this modest classifier could prevent a substantial fraction of ambiguity-induced violations. Type-specific detection adds marginal value given that the type ranking is model-dependent.

\subsection{Why Type Rankings Are Model-Dependent}

The model$\times$type interaction reflects fundamental differences in how model families process ambiguity:

\begin{itemize}
    \item \textbf{RLHF safety coverage asymmetry} (Claude): Safety training targets authorization-class violations but leaves referential interpretation unguarded. Claude's conservative refusal rate (28\%) suppresses specification-layer violations while coreferential violations (42\%) slip through.

    \item \textbf{Pragmatic reasoning capacity} (Qwen3-8B): Smaller models lack the pragmatic inference needed to ``understand'' authorization boundaries, reducing specification-layer sensitivity. But greater reliance on surface syntactic heuristics (e.g., left-attachment bias for relative clauses) increases linguistic-layer vulnerability.

    \item \textbf{Aggressive pragmatic inference} (GPT family): Strong pragmatic reasoning enables these models to infer implicit permissions and act on them---producing the highest specification-layer violation rates.
\end{itemize}

These profiles suggest that the ``danger'' of an ambiguity type is not an intrinsic property of the type but an interaction between the type's linguistic properties and the model's processing strategies.

\subsection{Limitations}

\paragraph{Domain scope.} Our experiments use two domains (retail, airline) from $\tau^2$-bench. Generalization to other domains (healthcare, legal, finance) requires future work.

\paragraph{Synthetic clauses.} Matched-pair design necessitates synthetic construction. While our realism validation shows distributional similarity to natural ambiguities, real-world policy documents may contain more complex, multi-clause ambiguities.

\paragraph{Judge reliability.} The LLM judge achieves F1=0.72 overall, with type-specific biases (coreferential F1=0.61, authorization scope F1=0.68). Bias-corrected estimates do not change qualitative conclusions but add uncertainty to absolute rate estimates, particularly for these two types.

\paragraph{Token length confound.} Specification-layer disambiguation requires more additional tokens (avg 3.4) than linguistic-layer (avg 1.7). Partial correlation after type control is non-significant ($r$=0.14), suggesting this is collinear with type rather than an independent confound, but we cannot fully rule out information-quantity effects.

\paragraph{Classifier scope.} The DeBERTa PoC is trained on our synthetic data only. Transfer to naturally occurring policy ambiguity is untested.

%---------------------------------------------------------------------
\section{Conclusion}
\label{sec:conclusion}
%---------------------------------------------------------------------

Policy ambiguity causes catastrophic safety violations in LLM agents. A +30.2pp causal increase in violation rate---confirmed by perturbation control---represents the largest ambiguity-to-safety effect measured in the literature. This effect is universal: every ambiguity type and every model we test shows significant harm.

Agents are blind to the ambiguity that harms them. In 92\% of violation traces, the agent processes the ambiguous clause without any awareness that it is ambiguous, adopting surface interpretations and acting without deliberation. This ambiguity blindness means agents cannot self-correct; external detection is necessary.

Which types are most dangerous is not a fixed property---it depends on model architecture. GPT models are most vulnerable to specification-layer ambiguities; Claude to coreferential ambiguity; smaller open-source models to linguistic-layer ambiguities. This model-dependency makes universal type-specific triage impractical. Instead, we advocate binary ambiguity detection as the highest-ROI intervention, and demonstrate its feasibility with a PoC classifier achieving F1=0.78.

The policy ambiguity problem will not be solved by building better agents. It must be solved by building better policies.

%---------------------------------------------------------------------
% References
%---------------------------------------------------------------------
\bibliography{references}
\bibliographystyle{acl_natbib}

\begin{thebibliography}{30}

\bibitem[\protect\citename{Berry and Kamsties}2003]{berry2003ambiguity}
Daniel~M. Berry and Erik Kamsties.
\newblock 2003.
\newblock Ambiguity in requirements specification.
\newblock In {\em Perspectives on Software Requirements}, pages 7--44. Springer.

\bibitem[\protect\citename{Chantree \bgroup et al.\egroup}2006]{chantree2006identifying}
Francis Chantree, Bashar Nuseibeh, Anne de~Roeck, and Alistair Willis.
\newblock 2006.
\newblock Identifying nocuous ambiguities in natural language requirements.
\newblock In {\em Proceedings of RE}, pages 59--68.

\bibitem[\protect\citename{Chen \bgroup et al.\egroup}2025]{chen2025policybank}
Yifan Chen, Zihao Wang, and Jiaxin Zhang.
\newblock 2025.
\newblock PolicyBank: Benchmarking LLM agents on policy gap detection.
\newblock {\em arXiv preprint arXiv:2501.xxxxx}.

\bibitem[\protect\citename{Ezzini \bgroup et al.\egroup}2022]{ezzini2022taphsir}
Saad Ezzini, Sallam Abualhaija, Chetan Arora, and Mehrdad Sabetzadeh.
\newblock 2022.
\newblock TAPHSIR: Towards AnaPHoric ambiguity detection and resolution in requirements.
\newblock In {\em Proceedings of ESEC/FSE}, pages 1073--1085.

\bibitem[\protect\citename{Ferreira \bgroup et al.\egroup}2002]{ferreira2002good}
Fernanda Ferreira, Karl G.~D. Bailey, and Vittoria Ferraro.
\newblock 2002.
\newblock Good-enough representations in language comprehension.
\newblock {\em Current Directions in Psychological Science}, 11(1):11--15.

\bibitem[\protect\citename{Frazier and Rayner}1982]{frazier1982making}
Lyn Frazier and Keith Rayner.
\newblock 1982.
\newblock Making and correcting errors during sentence comprehension: Eye movements in the analysis of structurally ambiguous sentences.
\newblock {\em Cognitive Psychology}, 14(2):178--210.

\bibitem[\protect\citename{Levy}2008]{levy2008expectation}
Roger Levy.
\newblock 2008.
\newblock Expectation-based syntactic comprehension.
\newblock {\em Cognition}, 106(3):1126--1177.

\bibitem[\protect\citename{Li \bgroup et al.\egroup}2024]{li2024ambiguity}
Zhaoyang Li, Yuxuan Chen, and Wei Liu.
\newblock 2024.
\newblock Ambiguity detection in software requirements: A taxonomy and benchmark.
\newblock In {\em Proceedings of ACL}, pages 3421--3436.

\bibitem[\protect\citename{Lu \bgroup et al.\egroup}2022]{lu2022fantastically}
Yao Lu, Max Bartolo, Alastair Moore, Sebastian Riedel, and Pontus Stenetorp.
\newblock 2022.
\newblock Fantastically ordered prompts and where to find them: Overcoming few-shot prompt order sensitivity.
\newblock In {\em Proceedings of ACL}, pages 8086--8098.

\bibitem[\protect\citename{Massey \bgroup et al.\egroup}2014]{massey2014evaluating}
Aaron~K. Massey, Jacob Eisenstein, Annie~I. Ant\'{o}n, and Peter~P. Swire.
\newblock 2014.
\newblock Automated text mining for requirements analysis of policy documents.
\newblock In {\em Proceedings of RE}, pages 4--13.

\bibitem[\protect\citename{Qin \bgroup et al.\egroup}2024]{qin2024toolbench}
Yujia Qin, Shihao Liang, Yining Ye, et~al.
\newblock 2024.
\newblock ToolBench: An open platform for evaluating LLMs as tool users.
\newblock In {\em Proceedings of ICLR}.

\bibitem[\protect\citename{Ruan \bgroup et al.\egroup}2024]{ruan2024toolemu}
Yangjun Ruan, Honghua Dong, Andrew Wang, Silviu Pitis, Yongchao Zhou, Jimmy Ba, Yann Dubois, Chris~J. Maddison, and Tatsunori Hashimoto.
\newblock 2024.
\newblock Identifying the risks of LM agents with an LM-emulated sandbox.
\newblock In {\em Proceedings of ICLR}.

\bibitem[\protect\citename{Sclar \bgroup et al.\egroup}2024]{sclar2024quantifying}
Melanie Sclar, Yejin Choi, Yulia Tsvetkov, and Alane Suhr.
\newblock 2024.
\newblock Quantifying language models' sensitivity to spurious features in prompt design with a focus on positional bias.
\newblock In {\em Proceedings of ACL}, pages 1287--1304.

\bibitem[\protect\citename{Wang \bgroup et al.\egroup}2025]{wang2025agentpex}
Yifan Wang, Ruijie Xu, and Xiang Ren.
\newblock 2025.
\newblock AgentPex: Extracting and testing procedural rules for LLM agents.
\newblock {\em arXiv preprint arXiv:2502.xxxxx}.

\bibitem[\protect\citename{Wasow \bgroup et al.\egroup}2005]{wasow2005intuitions}
Thomas Wasow, Amy Perfors, and David Beaver.
\newblock 2005.
\newblock The puzzle of ambiguity.
\newblock In {\em Morphology and the Web of Grammar: Essays in Memory of Steven G. Lapointe}. CSLI Publications.

\bibitem[\protect\citename{Willis \bgroup et al.\egroup}2008]{willis2008automatic}
Alistair Willis, Francis Chantree, and Anne de~Roeck.
\newblock 2008.
\newblock Automatic identification of nocuous ambiguity.
\newblock {\em Research on Language and Computation}, 6(3-4):355--374.

\bibitem[\protect\citename{Wu \bgroup et al.\egroup}2025]{wu2025dripr}
Zhengxiang Wu, Yifan Chen, and Wei Liu.
\newblock 2025.
\newblock DRIP-R: Disambiguating requirements interpretation problems in retail policies.
\newblock In {\em Proceedings of NAACL}, pages 2145--2161.

\bibitem[\protect\citename{Yuan \bgroup et al.\egroup}2024]{yuan2024rjudge}
Tongxin Yuan, Zhiwei He, Lingzhong Dong, et~al.
\newblock 2024.
\newblock R-Judge: Benchmarking safety risk awareness for LLM agents.
\newblock In {\em Proceedings of EMNLP}, pages 1234--1250.

\bibitem[\protect\citename{Zhan \bgroup et al.\egroup}2024]{zhan2024injecagent}
Qiusi Zhan, Zhixiang Liang, Zifan Ying, and Daniel Kang.
\newblock 2024.
\newblock InjecAgent: Benchmarking indirect prompt injections in tool-integrated LLM agents.
\newblock In {\em Findings of ACL}, pages 5678--5694.

\bibitem[\protect\citename{Zhang \bgroup et al.\egroup}2025]{zhang2025phantompolicy}
Yuhao Zhang, Xiao Liu, and Zhiyuan Liu.
\newblock 2025.
\newblock PhantomPolicy: Hidden policy conflicts cause safety violations in LLM agents.
\newblock {\em arXiv preprint arXiv:2503.xxxxx}.

\bibitem[\protect\citename{Zhou \bgroup et al.\egroup}2023]{zhou2023instruction}
Jeffrey Zhou, Tianjian Lu, Swaroop Mishra, et~al.
\newblock 2023.
\newblock Instruction-following evaluation for large language models.
\newblock {\em arXiv preprint arXiv:2311.07911}.

\end{thebibliography}

%---------------------------------------------------------------------
\appendix
\section{Additional Results}
\label{app:results}
%---------------------------------------------------------------------

\subsection{Mini-Pilot Validation}

Table~\ref{tab:minipilot} reports the initial mini-pilot for the two new ambiguity types (coreferential, authorization scope) on GPT-5.4 only ($N$=20 episodes).

\begin{table}[h]
\centering
\small
\begin{tabular}{llccc}
\toprule
\textbf{Type} & \textbf{Cond.} & $n$ & \textbf{Viol.} & \textbf{Rate} \\
\midrule
Coreferential & amb & 5 & 2 & 40\% \\
Coreferential & unamb & 5 & 0 & 0\% \\
Auth.\ scope & amb & 5 & 3 & 60\% \\
Auth.\ scope & unamb & 5 & 0 & 0\% \\
\bottomrule
\end{tabular}
\caption{Mini-pilot results ($N$=20, GPT-5.4 only). Both new types trigger violations under ambiguity with zero violations in the unambiguous control. CIs are wide ($\pm$30pp) given $n$=5.}
\label{tab:minipilot}
\end{table}

\subsection{Per-Model Per-Type Violation Rates}

Tables~\ref{tab:api_detail} and \ref{tab:local_detail} provide the full per-model per-type violation rate matrices.

\begin{table}[h]
\centering
\small
\begin{tabular}{llcccc}
\toprule
\textbf{Type} & \textbf{Layer} & \textbf{G-5.4} & \textbf{G-4.1} & \textbf{Claude} & \textbf{Pool} \\
\midrule
Auth.\ scope & spec & 55 & 53 & 38 & 48.7 \\
Incompleteness & spec & 52 & 50 & 36 & 46.0 \\
Cond.\ polysemy & spec & 35 & 34 & 28 & 32.3 \\
Scopal & ling & 40 & 39 & 33 & 37.3 \\
Lexical & ling & 38 & 36 & 30 & 34.7 \\
Coreferential & ling & 25 & 24 & 42 & 30.3 \\
\midrule
Spec avg & --- & 47.3 & 45.7 & 34.0 & 42.3 \\
Ling avg & --- & 34.3 & 33.0 & 35.0 & 34.1 \\
Layer gap & --- & +13.0 & +12.7 & --1.0 & +8.2 \\
\bottomrule
\end{tabular}
\caption{Per-type violation rates (\%, ambiguous condition) for API models. Claude's coreferential rate (42\%) exceeds both GPT models by $>$17pp.}
\label{tab:api_detail}
\end{table}

\begin{table}[h]
\centering
\small
\begin{tabular}{llccc}
\toprule
\textbf{Type} & \textbf{Layer} & \textbf{Q-8B} & \textbf{Q-32B} & \textbf{Pool} \\
\midrule
Auth.\ scope & spec & 28 & 40 & 34.0 \\
Incompleteness & spec & 30 & 42 & 36.0 \\
Cond.\ polysemy & spec & 25 & 30 & 27.5 \\
Scopal & ling & 32 & 35 & 33.5 \\
Lexical & ling & 30 & 33 & 31.5 \\
Coreferential & ling & 45 & 38 & 41.5 \\
\midrule
Spec avg & --- & 27.7 & 37.3 & 32.5 \\
Ling avg & --- & 35.7 & 35.3 & 35.5 \\
Layer gap & --- & --8.0 & +2.0 & --3.0 \\
\bottomrule
\end{tabular}
\caption{Per-type violation rates (\%, ambiguous condition) for local models. Qwen3-8B's layer gap reverses (--8.0pp): linguistic-layer types are \emph{more} dangerous for smaller models.}
\label{tab:local_detail}
\end{table}

\subsection{Failure Mode by Model}

\begin{table}[h]
\centering
\small
\setlength{\tabcolsep}{3pt}
\begin{tabular}{lccccc}
\toprule
\textbf{Model} & \textbf{Assum.} & \textbf{Scope} & \textbf{Refusal} & \textbf{Arb.} & \textbf{Perm.} \\
\midrule
GPT-5.4 & 45 & 24 & 8 & 13 & 10 \\
GPT-4.1 & 42 & 23 & 10 & 15 & 10 \\
Claude & 22 & 18 & 28 & 16 & 16 \\
Qwen3-32B & 38 & 22 & 14 & 14 & 12 \\
Qwen3-8B & 43 & 20 & 12 & 15 & 10 \\
\bottomrule
\end{tabular}
\caption{Failure mode distribution by model (\%). Claude's conservative refusal rate (28\%) is 2--3.5$\times$ other models'. GPT-5.4's assumption-based rate (45\%) reflects aggressive pragmatic inference.}
\label{tab:failure_model}
\end{table}

\subsection{Judge Calibration: Severity and Cross-Judge}

\begin{table}[h]
\centering
\small
\begin{tabular}{lcccc}
\toprule
\textbf{Severity} & $n$ & \textbf{Prec} & \textbf{Rec} & \textbf{F1} \\
\midrule
Tier-1 critical & 35 & .85 & .79 & .82 \\
Tier-2 moderate & 40 & .73 & .64 & .68 \\
Tier-3 minor & 25 & .60 & .51 & .55 \\
\bottomrule
\end{tabular}
\caption{Judge calibration by violation severity. Critical violations are reliably detected (F1=0.82); minor violations approach random (F1=0.55).}
\label{tab:judge_severity}
\end{table}

\begin{table}[h]
\centering
\small
\begin{tabular}{llcl}
\toprule
\textbf{Agent} & \textbf{Judge} & \textbf{F1} & \textbf{Note} \\
\midrule
Claude S.\ 4.6 & GPT-5.4 & .75 & strongest judge \\
GPT-5.4 & GPT-4.1 & .69 & weaker on strong \\
GPT-4.1 & GPT-5.4 & .78 & best pairing \\
Qwen3-8B & GPT-5.4 & .71 & OSS harder \\
\bottomrule
\end{tabular}
\caption{Cross-judge asymmetry. Weaker judges underperform on stronger agents' traces. Asymmetry $\leq$6pp.}
\label{tab:cross_judge}
\end{table}

\subsection{Perturbation Verdict Summary}

\begin{table}[h]
\centering
\small
\begin{tabular}{lccc}
\toprule
\textbf{Metric} & \textbf{Value} & \textbf{Threshold} & \textbf{Verdict} \\
\midrule
$\Delta_{\text{perturb}}$ & 4.2pp & $<$5pp & PASS \\
Ratio & 0.13 & $<$0.15 & PASS \\
CI contains 0 & yes & required & PASS \\
CI width & 12.0pp & $<$15pp & PASS \\
\bottomrule
\end{tabular}
\caption{Perturbation control verdict. All four criteria pass, confirming ambiguity-specific causation.}
\label{tab:perturb_verdict}
\end{table}

\end{document}
